Das $p(V)$-Diagramm zeigt einen Kreisprozess einer festen Menge eines idealen Gases. Die Kurven von B nach C und von D nach A liegen auf Hyperbeln (gestrichelt). Die Achsen sind in Einheiten von $p_0$, $V_0$ und $T_0 = \frac{p_0V_0}{nR}$ beschriftet; damit ist $T/T_0 = (p/p_0)\cdot(V/V_0)$.

\begin{center}
\begin{tikzpicture}[x=0.95cm, y=0.95cm, >=Stealth, font=\small]
 \draw[black!12] (0,0) grid[xstep=1, ystep=1] (5.001,5.001);
 \draw[->] (0,0) -- (5.6,0) node[right] {$V$ in $V_0$};
 \draw[->] (0,0) -- (0,5.6) node[above] {$p$ in $p_0$};
 \node[below left] at (0,0) {\small 0};
 \draw (1,0.08) -- (1,-0.08) node[below] {\small 1};
 \draw (2,0.08) -- (2,-0.08) node[below] {\small 2};
 \draw (3,0.08) -- (3,-0.08) node[below] {\small 3};
 \draw (4,0.08) -- (4,-0.08) node[below] {\small 4};
 \draw (5,0.08) -- (5,-0.08) node[below] {\small 5};
 \draw (0.08,1) -- (-0.08,1) node[left] {\small 1};
 \draw (0.08,2) -- (-0.08,2) node[left] {\small 2};
 \draw (0.08,3) -- (-0.08,3) node[left] {\small 3};
 \draw (0.08,4) -- (-0.08,4) node[left] {\small 4};
 \draw (0.08,5) -- (-0.08,5) node[left] {\small 5};
 \draw[black!35, dashed, domain=1.6:5, samples=50, smooth] plot (\x, {8/\x});
 \draw[black!35, dashed, domain=0.8:5, samples=50, smooth] plot (\x, {4/\x});
 \draw[very thick, blue] (1,4) -- (2,4);
 \draw[-{Stealth[length=2.6mm]}, blue] (1.54,4) -- (1.56,4);
 \draw[very thick, blue, domain=2:4, samples=60, smooth] plot (\x, {8/\x});
 \draw[-{Stealth[length=2.6mm]}, blue] (2.821,2.835) -- (2.835,2.821);
 \draw[very thick, blue] (4,2) -- (4,1);
 \draw[-{Stealth[length=2.6mm]}, blue] (4,1.46) -- (4,1.44);
 \draw[very thick, blue, domain=1:4, samples=60, smooth] plot (\x, {4/\x});
 \draw[-{Stealth[length=2.6mm]}, blue] (2.007,1.993) -- (1.993,2.007);
 \fill[black] (1,4) circle (2pt) node[above left] {$A$};
 \fill[black] (2,4) circle (2pt) node[above right] {$B$};
 \fill[black] (4,2) circle (2pt) node[above right] {$C$};
 \fill[black] (4,1) circle (2pt) node[below left] {$D$};
\end{tikzpicture}
\end{center}

\begin{abcliste}
\abc \DuSie{Benenne}{Benennen Sie} die vier Schritte.
\abc Welche Aussagen sind richtig, welche falsch? \DuSie{Begründe}{Begründen Sie}.
\begin{enumerate}
\item Von B nach C halbiert sich der Druck, während sich das Volumen verdoppelt.
\item Im Zustand C ist die Temperatur gleich hoch wie im Zustand A.
\item Von A nach B verdoppelt sich die Temperatur.
\item Die Hyperbel durch D und A gehört zur höheren Temperatur.
\item Von C nach D sinkt die Temperatur auf die Hälfte.
\item Von D nach A wird das Gas bei gleichbleibendem Druck komprimiert.
\end{enumerate}
\end{abcliste}
